\chapter{Discussion}
\label{chap:discussion}

\section{Two regimes, two answers}
\label{sec:discussion-regimes}
Our central empirical pattern is a contrast between regimes. In the high-diversity regime, on
ASL Citizen~\cite{aslcitizen}, the appearance interventions returned a null: no
objective family beats plain cross-entropy at matched depth, and the
frozen backbone's linear probe stays the best in
Table~\ref{tab:backbone_ft_retrieval} on every probe metric, so nothing we
trained made the task more accessible than the untouched representation
already made it. In the
low-diversity NGT744 setting, where the model must bridge a real signer gap from a
handful of recordings, motion-locked synthetic counterfactuals produced
a clear gain. The simplest reading is that appearance invariance is not a
property a model has or lacks in the abstract; it is relative to the
diversity already present in training. The Logos checkpoint we
use~\cite{ovodovLogosWellTemperedPretrain2025} is trained jointly on three
isolated-sign corpora (Logos, AUTSL, and WLASL); the Logos dataset alone
contributes roughly $200{,}000$ videos from $381$ signers, and AUTSL and
WLASL add $43$ and $119$ more.
\todome{REWRITTEN 12.08: the previous sentence attributed the contrast to
backbone pretraining exposure, but both regimes run on the same frozen
checkpoint, so pretraining cannot be the variable that separates them.
Signer counts moved to a consistent whole-corpus basis (was 305 train-split
Logos signers plus a wrong 36 for AUTSL; 43 per the ChaLearn description
in Related Work).}
Both regimes run on this same frozen checkpoint, so pretraining exposure
is a constant across the contrast, not the variable; what differs is the
training data each intervention enters. The ASL Citizen runs already see
$35$ real signers recorded under varied webcam conditions, so the edits
add appearance variation of a kind the split carries anyway. The NGT744
encoder trains on four signers in one studio, and there the synthetic
signers supply appearance diversity the training data otherwise lacks.

% Signer counts, whole-corpus basis (12.08): Logos paper §3.1 gives
% 199,668 videos / 381 signers; AUTSL 43 (ChaLearn, Related Work §2);
% WLASL 119. The paper prints no train-split signer integer, so the
% earlier ~305 train-only estimate was dropped.

A second factor compounds this signer-diversity gap for the raw, untrained
backbone specifically. That Logos transfers well onto ASL Citizen even
before any ASL-Citizen-specific training
(Table~\ref{tab:backbone_ft_retrieval}: the unmodified frozen backbone
already exceeds the benchmark's published baselines with no target-corpus
training at all)
makes sense once the pretraining recipe is considered: Logos is trained
jointly with WLASL, and WLASL and ASL Citizen share a large fraction of
their vocabulary: both are American Sign Language isolated-sign corpora
drawn from an overlapping sign inventory, so the backbone has effectively
already seen much of what ASL Citizen asks of it. The shift to NGT744 removes
that overlap on three axes at once: NGT744 is in a different sign
language entirely (Dutch Sign Language rather than American), its
recordings come from a different capture setup (signers in motion-capture
suits in a studio, rather than webcam video), and its unit of analysis is
sentence-level continuous signing rather than the isolated, gloss-level
clips both ASL Citizen and WLASL provide. Any one of these differences
would weaken transfer; together they are a plausible reason why the same
pretrained backbone that carries over to ASL Citizen with no
target-corpus training needs
the full synthetic-appearance intervention to bridge the NGT744 signer gap.

This reading also bears on SDDA's positive
result~\cite{fuSignerDiversitydrivenData2024b} on PHOENIX-2014T, with a
caveat that has to come first: SDDA reports sign language translation on
continuous video where we report isolated-sign retrieval, so what follows
compares regimes, not two measurements of the same quantity. Their
generatively augmented contrastive setup operates on a studio-recorded corpus
with far less signer variety than ASL Citizen (PHOENIX-2014T contains nine signers in
total~\cite{fuSignerDiversitydrivenData2024b}), i.e.\ closer to our NGT744
% Nine signers verified 2026-07-12 in three independent sources: Camgoz et
% al. 2018 §4 (primary; no bib entry yet — optional to add), the SDDA paper
% itself ("includes data from only nine different signers"), and the Isharah
% paper's comparison table. Cited to SDDA since it is already in the bib and
% is the paper whose corpus is being discussed.
regime than to our ASL Citizen regime.

In PHOENIX-2014T, recording conditions are held fixed (one studio, one camera
setup, one weather-forecast register), so the shift their augmentation has
to bridge is signer appearance and signing style alone, narrower than the
corpus and viewpoint changes our NGT744 evaluation layers on top of it. How
much the held-out-signer choice matters on this corpus is documented
independently: under a full signer-disjoint re-evaluation of
PHOENIX-2014T, mean translation scores fall by between a third and four
fifths depending on the system, and the choice of held-out signer moves
BLEU-4 by up to a factor of ${\sim}5$ for the most affected
one~\cite{artiagaRethinkingSignLanguage2025}.
\section{RGB versus pose, revisited}
\label{sec:discussion-rgb-pose}
Pose input is widely adopted in SLR partly \emph{because} it is assumed
to discard appearance: the robustness of pose-based models is
attributed to inherent invariance to clothing, background and
lighting~\cite{moryossefEvaluatingImmediateApplicability2021,
kratimenosIndependentSignLanguage2020, tranRegionAwarePoseModeling2025},
% keys as used at lines 423-424
and the SignEval organisers went as far as distributing
\emph{only} pose features for a signer-independent task, citing signer
overfitting and privacy~\cite{luqmanSignEval2025Challenge2025a,
alyamiIsharahLargeScaleMultiScene2026b}. % keys as used at lines 182-183
Our diagnostic probes complicate the premise behind both moves. On ASL
Citizen, mean-pooled MediaPipe keypoints are the \emph{most}
signer-separable of the three raw feature spaces we test, on both
readouts: their inter/intra distance ratio is
$29.9\times$ [$28.4$, $32.0$] % table line 1360 (tab:raw_feature_l2_probe)
against
$22.7\times$ and $22.2\times$ % lines 1361-1362 (Logos, I3D)
for the two RGB backbones (confidence intervals overlapping
neither's), % prose lines 1372-1374
and a linear head recovers the signer at
$98.7\%$ % line 1360 (MediaPipe decoding; chance 3.0%, caption line 1354
% and prose lines 1370-1371)
against $94.2\%$ (I3D) and $65.9\%$ (Logos), an ordering unchanged on
the signer-disjoint test split
($99.4\%$/$95.2\%$/$74.3\%$, chance $9.1\%$). % lines 1374-1376
One asymmetry bounds the reading: the two RGB spaces come from
gloss-supervised backbones, which \S\ref{sec:results-l2-raw} shows sheds
linearly accessible identity, while the keypoints are untrained
coordinates, so the ordering measures what pose \emph{retains}, not what
the modality would concentrate at matched training.
Camera framing cannot carry this: the keypoints are root-centred at the
shoulder midpoint and scale-normalised to shoulder width, so what
identifies the signer is shoulder-normalised geometry: body
proportions, habitual posture, movement style. % prose lines 1376-1381
Nor is it gloss content in disguise: under gloss-disjoint
cross-validation on ASL Citizen, where the classifier is tested only on
signs no signer contributed during its training, the decoding numbers
are essentially unchanged ($98.7\%$/$94.6\%$/$64.5\%$,
\S\ref{sec:results-l2-raw}).
Pose discards pixels; on this evidence it retains precisely the
person-specific signal the modality is assumed to remove.

This lands on the same side as the fairness literature reviewed in
\S\ref{sec:related-work-modality}, from the opposite direction. That literature's caveat is about
the \emph{extractor}: pose pipelines are themselves learned RGB models
whose training and evaluation sets under-represent darker-skinned
individuals~\cite{lachanceCaseStudyFairness2023,
hasanHi5SyntheticData2024}, % keys as used at line 194
whose audits reach conclusions that flip with the evaluation
set~\cite{lachanceCaseStudyFairness2023,
xiangFairHumanCentricImage2025}, % keys as used at line 203
whose hand tracking degrades exactly where signing is linguistically
dense~\cite{moryossefOptimizingHandRegion2024,
moryossefEvaluatingImmediateApplicability2021}, % keys as used at lines 218-221
and whose ASL Citizen pose baseline shows larger skin-tone and gender
disparities than its RGB
counterpart~\cite{atwellStudyingMitigatingBiases2024}. % key as used at line 210
Our probes add the complementary caveat about the \emph{representation}:
even when extraction succeeds and coordinates are normalised, the
resulting features are more signer-identifying than either RGB space we
compare against. Together the two caveats suggest that treating pose
distribution as an appearance-bias or privacy fix (as in the
pose-only challenge protocol above) deserves the same scrutiny as
any other debiasing claim: the appearance axis a pose skeleton removes
(clothing, background, skin pixels) is not the only axis along which
signers are identifiable, and by our measurements not the dominant one.

The scope of this claim should be stated as narrowly as the evidence.
Ours is a bias probe, not a matched-capacity recognition comparison: we
never train a pose-based recogniser ourselves (the pose rows of
Table~\ref{tab:asl_citizen_retrieval} are the published results of the
original SignCLIP model), % setup lines 1058-1061
so we make no claim about whether pose-based \emph{models} exploit this
identity signal, how much it costs them under signer shift, or how the
modalities compare at matched training: questions the recent
RGB-versus-pose results
cited in \S\ref{sec:related-work-modality}~\cite{minCloserLookSkeletonbased2025,
wongSignRepEnhancingSelfSupervised2025} % keys as used at lines 419, 240
address from the accuracy side. 
What the probes do establish is the narrower, assumption-level point:
signer identity is not removed by pose extraction by construction, and
a pipeline that adopts pose \emph{for} invariance or privacy inherits
an unexamined assumption rather than a guarantee.

\todome{REWRITTEN 12.08: the previous version compared the heads'
train-split $\rho$ against the test-split reference ($17.67\times$),
which the table caption itself forbids; against the train-split frozen
ratio the geometric half of the claim disappears. The matching caption
fix resolves the `framing check' todo that sat on
tab:head\_signer\_probes.}
After training, the two readouts come apart. The SignCLIP-style heads
make identity far \emph{less} linearly decodable ($20$--$25\%$ against
$66.9\%$ for the frozen features on the same split, chance $2.9\%$)
while leaving the distance geometry near the frozen level
($22.0$--$23.4\times$ against $22.7\times$,
Table~\ref{tab:raw_feature_l2_probe}), with only the $K{=}5$ pair-SupCon
head clearly above it ($26.3\times$,
Table~\ref{tab:head_signer_probes}). By the bar set in
\S\ref{sec:results-ft-probes}, that an intervention which genuinely
removes identity must lower both readouts, no objective in this thesis
removed signer structure; training moved where that structure is
measurable, which is the reason to report both probes rather than treat
linear decodability as \emph{the} appearance-bias metric.

\section{Limitations}
\label{sec:limitations}

\textbf{Untrimmed NGT744 recordings.} The NGT744 training videos were not
trimmed: each begins with a lengthy idle segment in which the signer
stands watching the reference video before starting to sign, and ends
with a reach toward the foot pedal used to stop the recording. Both add
frames unrelated to signing (idle standing and extraneous motion),
and because the idle portion differs in length per signer, no trivial
automated trim existed; we judged manual cleaning not worth thesis time.

\textbf{Frame-independent generative editing introduces flicker.} Our
editing pipeline (\S\ref{subsection:GAE}) edits each frame of a
source video independently: every edited frame follows the appearance-edit
prompt on its own, with no explicit mechanism enforcing consistency with
the frame before or after it, so consecutive edited frames drift slightly
relative to one another and the resulting clip flickers. We looked for a
video-native alternative that would edit a clip as one temporally
consistent unit rather than frame-by-frame, trying LTX 2.3 (22B)
~\cite{hacohen2026ltx2}, LTX 2.3 with an Edit Anything
LoRA~\cite{alissonerdx2026editanything}, Wan
2.2~\cite{wan2025}, and Lucy Edit 1.1 dev
(5B)~\cite{decart2025lucyedit};
in every case the edited
output was, on visual inspection, substantially worse than the
frame-by-frame FLUX.2-klein-9B route, so we
kept the per-frame pipeline and report the resulting flicker as a
limitation rather than a solved problem.

\textbf{Occasional generative hallucinations.} The editing model
occasionally hallucinates content that was not in the source frame:
most visibly, adding a spurious or malformed handshape to a frame whose
original signer was not making that shape, or to a frame that did not
contain a visible hand at all. This risk is highest for the skin-tone and
signer-swap edit types (\S\ref{subsection:GAE}), and is a direct
consequence of editing each frame independently rather than conditioning
on the sign actually being performed.
 A more heavily
constrained editing pipeline (for instance, conditioning the editor
on the source pose via ControlNet-style
conditioning~\cite{zhangAddingConditionalControl2023}) could reduce these hallucinations
further; we did not quantify their prevalence, and judged the added
pipeline complexity out of scope for this thesis.

\todome{ADDED 12.08: experiment-level limitations; the previous three
items covered only the generation tooling. Cut whatever you disagree
with.}
\textbf{Statistical and coverage limits of the ASL Citizen experiments.}
Every ASL Citizen cell is a single training run, without seeds or
confidence intervals, against a nuisance floor that is not zero (probe
retest spread up to $1.8$ points; a $1.8$-point retrieval shift from a
change of frame preprocessing alone); the NGT744 experiments carry three
seeds and per-seed significance tests, the ASL Citizen ones do not. The
generative edits reach $450$ of the $2{,}731$ glosses ($14.9\%$ of
training videos), and the corpus's minimum of eight signers per gloss
bounds the headroom any signer-diversity intervention can have there
(\S\ref{sec:setup-generation}). Finally, the two regimes differ in more
than signer diversity: sign language, capture setup and unit of analysis
all change between ASL Citizen and NGT744
(\S\ref{sec:discussion-regimes}), and the from-scratch control that
would isolate the diversity account remains future work
(\S\ref{sec:future-work}).

\section{Future work}
\label{sec:future-work}
\textbf{Recovering motion from video rather than from motion capture.}
The motion-locked route requires motion-capture recordings, which is why
it could not be applied to a corpus that provides only RGB video
(\S\ref{sec:method-diffusion}). Estimating parametric
SMPL-X~\cite{pavlakosExpressiveBodyCapture2019a} pose and shape from the
video itself would replace that input, so that the recovered parameters
drive the same rendering and the route extends to any corpus.
SignAvatars~\cite{yuSignAvatarsLargeScale3D2025} and Neural Sign
Actors~\cite{baltatzisNeuralSignActors2024a} show what such
parametrisation looks like for signing, though both would need to improve
substantially before they could stand in for recorded capture.

\textbf{Signer identity persists through the synthetic-avatar swap.}
Even
when signer identity is changed by re-rendering with a synthetic avatar,
the avatar does not fully erase who provided the underlying motion
capture: body shape, body posture, range of motion when signing, and
signing speed are all retargeted along with the sign itself, and all four
remain visibly present in the render. This is consistent with the
diagnosis that shoulder-normalised body geometry and movement style
identify signers (\S\ref{sec:results-diagnosis}), and it is plainly
visible: an observer with no knowledge of sign language can still tell,
from individual frames of the renders, which real NGT744 signer supplied
the motion capture for a given avatar. The synthetic-signer
intervention changes surface appearance (body model, clothing, skin
tone) without fully removing the motion-level signature of the source
signer, a distinction the retrieval-based evaluation in this thesis does
not separate out. Measuring that motion-level signature directly (a
signer-identity probe on the rendered views themselves) would be a
cheap first step, and motion retargeting that also perturbs kinematic
style the harder follow-up.

\textbf{Augmenting the full ASL Citizen training set, not just a targeted
subset.} Our generative augmentation covers a deliberately targeted subset
of the ASL Citizen training set, chosen worst-first by signer scarcity and
baseline difficulty (\S\ref{sec:method-diffusion}). The subset itself was
set by the generation compute budget (\S\ref{sec:setup-generation});
evidence that targeted synthetic augmentation can outperform augmenting
everything~\cite{nguyenDoWeNeedSynthetic2025} guided how the subset was
chosen, not whether to use one.
Whether that finding transfers to our setting is untested: it remains an
open question whether augmenting the entire training set, rather than the
selected worst-first subset, would have produced a larger gain
than the targeted budget did, or whether the targeted-subset advantage
holds here too and full-corpus augmentation would mostly add compute cost
without adding retrieval gain.

\textbf{End-to-end training for the NGT744 encoder.} All NGT744 results come
from an encoder trained on pre-extracted, frozen Logos clip features
(\S\ref{sec:results-lowdiv}). An end-to-end variant that prepends
the MViTv2-S backbone itself (a frozen-backbone control and a fully
fine-tuned configuration) was implemented and configured but never
run. Running both would show whether the low-diversity gain from
synthetic views survives, grows, or shrinks once gradients can reach the
visual backbone, giving the low-data regime its own counterpart to the
contrast between frozen-feature heads and backbone
fine-tuning (\S\ref{sec:discussion-regimes}).

\textbf{A discriminative term for the high-$K$ arms.} Packing more
synthetic views into a training item does not pay: on NGT744, stacking
side-view renders onto the matched arm drops R@1 from $0.649$ to $0.437$
(E3 to E5), adding further generative edits drifts it from $0.505$ to
$0.461$ (E4 to E7), and the pair-SupCon head points the same way on ASL
Citizen ($0.567$ at $K{=}3$ falling to $0.414$ at $K{=}5$). Schedules are
matched to each arm's update count and the contrastive batch is fixed at
$64$ items for every arm (\S\ref{sec:setup-training}), so neither
truncated training nor negative starvation accounts for the drops. Our
reading is that pure SupCon at high $K$ spends its capacity making views
of the same item indistinguishable, at the cost of keeping different
items apart. A
natural test is to add a discriminative term to the NGT744 training setup
(a 12-layer encoder trained with SupCon on frozen Logos clip features;
\S\ref{sec:results-lowdiv}): cross-entropy on sentence-recording
identity plus $\lambda$ times the SupCon term, sweeping
$\lambda \in \{0.1, 1\}$ over a subset of arms. If the high-$K$ drop
disappears, the tradeoff is objective-specific and fixable; if it
persists, appearance invariance itself is what costs retrieval.

\textbf{A from-scratch control for the high-diversity null.} The two-regime
reading attributes the flat high-diversity result to appearance
invariance already supplied by real diversity, in pretraining and in the
35-signer training split (\S\ref{sec:discussion-regimes}), and the
pretraining half of this
is directly testable: train the same MViTv2-S from random initialisation
on ASL Citizen, with cross-entropy alone and with the synthetic
augmentations added, and compare against the pretrained baseline. If the
from-scratch model lands well below the pretrained one and the
augmentations now help it, pretraining diversity, not a saturated
recipe, explains the null; if augmentation still adds nothing, the
ceiling is not a pretraining artifact. 

\textbf{Fine-tuning I3D as a second backbone for the embedding-space
analyses.} Throughout this thesis I3D is used only as a frozen,
off-the-shelf feature extractor (\S\ref{sec:results-diagnosis},
\S\ref{sec:results-highdiv}), never fine-tuned the way the Logos/MViTv2-S
backbone is. Repeating the fine-tuning runs and the embedding-space
analyses (signer-separability probes, form-dose geometry) with a
fine-tuned I3D would show whether this thesis's diagnostic findings are
specific to the MViTv2-S/Logos backbone or hold across RGB architectures
more generally.
